\section{Comparison with Plain Entailment}
\label{sec:comparison}

Plain entailment $\Gamma\entails x\le y$ is the case $k=1$, and its decision procedure
\cite{NSSE26} also bounds the shortest holes of readers. The objects are the same: readers,
holes, comparisons, word actions, and the shortening of blocks on a reverse trie. The
argument for one word does not extend directly to several, and each step of our proof
replaces one step of the one-word argument. \Cref{tab:comparison} summarizes the
correspondence.

\begin{table}[!htb]
\centering\small
\caption{The bound on shortest holes for one word \cite{NSSE26} and for a tuple of words.}
\label{tab:comparison}
\begin{tabular}{@{}>{\raggedright\arraybackslash}p{0.2\textwidth}>{\raggedright\arraybackslash}p{0.33\textwidth}>{\raggedright\arraybackslash}p{0.39\textwidth}@{}}
\toprule
 & One word & A tuple of words (this paper)\\
\midrule
comparisons & within the word & within and between the words\\
after the first rejection & no admitted comparison & a cut can still be a host
  (\cref{fig:stopping})\\
finite holes & end at a dead end & stopped words: every letter dies or is blocked by an
  earlier stopped word (\cref{lem:stopping})\\
infinite holes & an eventually periodic word & rays with bounded periods
  (\cref{lem:periodic})\\
depth & bounded at a dead end & bounded on the stopped words by charging
  (\cref{lem:depth})\\
shortening & one word, on a reverse trie & all words together, linked blocks, buffers
  (\cref{lem:refinement,lem:buffered})\\
cutting & a periodic word is cut directly & strict difference constraints
  (\cref{lem:truncation})\\
\bottomrule
\end{tabular}
\end{table}

\subparagraph*{Rejection no longer ends the comparisons.}
For one word, the prefix test puts the host cut at or before the target cut, and a target
cut after the first rejection carries a live successor of a rejected state, which is never
admitted. An admitted comparison therefore has both cuts at or before the first rejection,
and a hole either ends at a \emph{dead end} or extends through rejected states forever, and
then periodically. A dead end is a rejected state from which every letter leads to a state
that is not rejected. In a structural reader this means that every letter leads to the
dead state. With several words the host and the target lie in different words. A cut far
beyond the first rejection of one word can be the host of a comparison whose target lies
early in another word (\cref{fig:stopping}). In \cref{ex:comparable-reader} the first word
can neither be cut at its first rejected cut nor be extended by the letter $1$. Stopping
(\cref{sec:stopping}) is the replacement for dead ends. A stopped word is a word whose
every letter dies, as at a dead end, or is blocked by a word that stopped earlier. Only
fixed words can block: when the target grows together with the host, a comparison at the
new end of the target has a live successor of a rejected state as its target state, which
is never admitted (\cref{lem:joint}).

\subparagraph*{Depth at stopped words.}
For one word, a hole that ends at a dead end has depth at most $2v+1$ for $v$ variables
\cite{NSSE26}. Whichever letter follows the dead end, the constraints force a leaf, and the
argument charges the nested recurrences of the end of the word to distinct variables. Here
the
suffixes of all stopped words nest along one host, and the bound charges each nesting to
the obstruction that stops it: the dead state, or an earlier stopped word, to which the
charge is transferred (\cref{lem:depth}). The charges give at most $4v+1$ lengths per
stopped word for a compiled reader with $v$ variables.

\subparagraph*{Periodic parts must survive the shortening.}
For one word, a hole that extends forever through rejected states has an eventually
periodic extension; after its beginning is replaced by a short word, a prefix of quadratic
length is already a hole. Here a moving word can be
the host of comparisons with the stopped words and with the other rays, at positions far
beyond its first rejection, so its periodic part cannot simply be cut. The shortening keeps
a periodic buffer after each stem (\cref{sec:buffers}). The buffer keeps the comparisons
with stopped words visible in a bounded prefix, and keeps the relative offsets that
separate admitted comparisons between rays.

\subparagraph*{One trie for all words.}
For one word, the blocks of a level are shortened together on the trie of their reversals,
so that equal blocks stay equal and different blocks keep differing at the same place. With
several words, deleting a block from one word can create a comparison with another. The
markers are therefore the last letters of \emph{all} words, the blocks carry their front
markers so that a shortened list remains a parse of a word, and one dictionary of contents
is shortened for all coordinates at once (\cref{sec:trie}).

\subparagraph*{Cutting the rays.}
Finally, the rays have to be cut at finite lengths. For one word, no two cuts of the
periodic word with equal infinite rests carry an admitted pair, since every long prefix is a
hole, and a single cut a fixed distance after the first rejection suffices. Rays that grow
in step can have such pairs of cuts, between different rays or within one ray, as long as
the target rest stays longer than the host rest at every common increment. Two cuts with
equal infinite rests form a comparison exactly when the rest of the target is not longer
than the rest of the host, so the cutting lengths must satisfy strict inequalities, a
system of difference constraints with a bounded solution (\cref{lem:truncation}).

Both bounds are singly exponential. For one word, the exponent is a polynomial of smaller
degree in the number of variables \cite{NSSE26}; here it is $(4v+k+16)^{10}$. The
polynomial in the exponent is larger, because the depth of the stopped words grows with the
number of queries, and the buffers and the trie descent run over all words at once.
